The reduction at the $S_{2}$ point: for $\sqrt{m}$-triple $(\sqrt{t}, \sqrt{t}, \sqrt{T})$, $Q = (2t + T)/(2\sqrt{t} + \sqrt{T})^{2} = 2/3$ rearranges to $\rho^{2} - 8\rho - 2 = 0$ for $\rho = \sqrt{T/t}$, whose positive root is $4 + \sqrt{18} = 4 + 3\sqrt{2}$; in the Fourier coordinates of Proposition \claimref{families.chart-computation-sign-twist-bracket-form}(3) the $S_{2}$ point is $\delta = 0$ with $\sqrt{m} = (a + |\Gamma|, a - |\Gamma|/2, a - |\Gamma|/2)$, and at $|\Gamma| = a\sqrt{2}$ the ratio $(a + |\Gamma|)/(a - |\Gamma|/2) = (1 + \sqrt{2})(2 + \sqrt{2}) = 4 + 3\sqrt{2}$, the same root; the fourth power evaluates in $\mathbb{Z}[\sqrt{2}]$ as $(4+3\sqrt{2})^{2} = 34 + 24\sqrt{2}$ and $(34+24\sqrt{2})^{2} = 2308 + 1632\sqrt{2}$. \eqref{eq:crossing} is the Koide condition read on the skeleton spectrum of \claimref{families.skeleton-rank-one-plus-ordering-shift}, whose two light eigenvalues coincide; its middle form is the division $6x^{2} + 6x + 1 = 6x(1+x) + 1$; it is strictly increasing in $x > 0$ (derivative $6 - 1/(1+x)^{2} > 0$), which gives the single crossing and the displayed identity for $x^{*}$. The script brackets $x^{*}$ between rational bounds by exact sign evaluation in $\mathbb{Z}[\sqrt{2}]$ arithmetic, $769.332 < x^{*} < 769.333$, and confirms $K/6 - x^{*} \in (2 \times 10^{-4}, 3 \times 10^{-4})$. The conclusion clause is the paper's pre-registered refutation condition for a canonical depth, applied: the nominated canonical depth fails by orders of magnitude ($\tanh^{2}r = \tfrac12$, that is $x = 1$, gives a ratio of $13/2$), and no independent structure marks $x^{*}$.
